% ============================================================================
%  PRECÁLCULO — Sesión 5
%  Semana 5: 7 al 12 de septiembre de 2026
%  1.5 Razones y proporciones · Repaso integrador de la Unidad 1
%  Preparación del Examen 1
%  Duración prevista: 3 horas
% ============================================================================
\documentclass[aspectratio=169,10pt]{beamer}

\usepackage[utf8]{inputenc}
\IfFileExists{spanish.ldf}{\usepackage[spanish,es-noshorthands]{babel}}{}
\usepackage{amsmath,amssymb}
\usepackage{booktabs}
\usepackage{tikz}

\usetheme{Uniminuto}

\uniasignatura{Precálculo}
\uniprofesor{Andrés Felipe Puerta Velez}
\unicorreo{andres.puerta-v@uniminuto.edu.co}
\uninrc{MATE1161}
\unisemana{Sesión 5 — Semana 5 · 7 al 12 de septiembre de 2026}

\newcommand{\listatiny}{\setbeamerfont{itemize/enumerate body}{size=\tiny}\tiny}

\begin{document}

\portada

\begin{frame}{Orden del día}
  \begin{enumerate}
    \item Llamado a lista y revisión de los ejercicios propuestos
    \item \azul{Tema 1.5: razones, proporciones y regla de tres}
    \item Proporcionalidad directa e inversa; regla de tres compuesta
    \item Porcentajes, repartos proporcionales y escalas
    \item Repaso integrador de la Unidad 1
    \item Simulacro del Examen 1 y recomendaciones
  \end{enumerate}

  \vspace{0.25cm}
  {\footnotesize\color{uniGris} Sesión de 3 horas. La próxima semana
  \destacar{no hay clase nueva}: se presenta el Examen 1 (35\%) en Aulas
  Virtuales, presencial en el aula.}
\end{frame}

\seccion{1.5 Razones y proporciones}

\begin{frame}{Razón y proporción}
  \scriptsize
  \begin{columns}[T,onlytextwidth]
    \begin{column}{0.46\textwidth}
      \begin{block}{Razón}
        \tiny
        Comparación de dos cantidades por cociente:
        \[ a : b \quad\text{o}\quad \frac{a}{b} \]
        ``La razón de hombres a mujeres es $3:5$'' significa que por cada 3
        hombres hay 5 mujeres.
      \end{block}
    \end{column}
    \begin{column}{0.50\textwidth}
      \begin{block}{Proporción}
        \tiny
        Igualdad de dos razones:
        \[ \frac{a}{b} = \frac{c}{d} \quad\Longleftrightarrow\quad
           \destacar{a\cdot d = b\cdot c} \]
        Los \azul{extremos} son $a$ y $d$; los \azul{medios}, $b$ y $c$.
      \end{block}
    \end{column}
  \end{columns}

  \vspace{0.1cm}
  \begin{exampleblock}{Cuidado: razón no es lo mismo que fracción del total}
    \scriptsize
    Si la razón es $3:5$ y hay 240 personas, \destacar{no} son 3/5 hombres.
    Hay $3+5 = 8$ partes:
    \[ \text{una parte} = \frac{240}{8} = 30 \quad\Rightarrow\quad
       \text{hombres} = 3(30) = \destacar{90}, \quad
       \text{mujeres} = 5(30) = \destacar{150} \]
  \end{exampleblock}

  \vspace{0.05cm}
  {\tiny\color{uniGris} \azul{Verificación:} $90+150 = 240$ \checkmark \ y \
  $90:150 = 3:5$ \checkmark. \destacar{Toda respuesta de proporciones se puede
  verificar}: hágalo siempre.}
\end{frame}

\begin{frame}{Hallar el término desconocido}
  \scriptsize
  \begin{columns}[T,onlytextwidth]
    \begin{column}{0.48\textwidth}
      \begin{block}{Cuarta proporcional}
        \tiny
        \[ \frac{4}{6} = \frac{x}{15} \]
        Productos cruzados: $6x = 4(15) = 60$
        \[ x = \destacar{10} \]
        Verificación: $\tfrac{4}{6} = \tfrac{10}{15} = \tfrac{2}{3}$ \checkmark
      \end{block}
    \end{column}
    \begin{column}{0.48\textwidth}
      \begin{block}{Media proporcional}
        \tiny
        Cuando los medios son iguales:
        \[ \frac{4}{x} = \frac{x}{9} \;\Rightarrow\; x^2 = 36
           \;\Rightarrow\; x = \destacar{6} \]
        Es la \azul{media geométrica}: $x = \sqrt{a\cdot b}$.
      \end{block}
    \end{column}
  \end{columns}

  \vspace{0.08cm}
  \begin{alertblock}{Propiedades útiles de las proporciones}
    \tiny
    Si $\dfrac{a}{b} = \dfrac{c}{d}$, entonces también
    $\dfrac{b}{a} = \dfrac{d}{c}$, \ $\dfrac{a}{c} = \dfrac{b}{d}$ \ y \
    $\dfrac{a+b}{b} = \dfrac{c+d}{d}$. La última se usa cuando el enunciado da
    el \azul{total} en lugar de una de las partes.
  \end{alertblock}
\end{frame}

\begin{frame}{Proporcionalidad directa e inversa}
  \scriptsize
  \begin{columns}[T,onlytextwidth]
    \begin{column}{0.48\textwidth}
      \begin{block}{Directa}
        \tiny
        Si una sube, la otra sube en la misma proporción. El
        \azul{cociente} es constante:
        \[ \frac{y}{x} = k \qquad y = kx \]
        \textbf{Ejemplo:} 5 servidores procesan 1\,200 peticiones/s.
        ¿8 servidores?
        \[ \frac{1\,200}{5} = \frac{x}{8} \Rightarrow
           x = \destacar{1\,920} \]
      \end{block}
    \end{column}
    \begin{column}{0.48\textwidth}
      \begin{alertblock}{Inversa}
        \tiny
        Si una sube, la otra baja. El \azul{producto} es constante:
        \[ x\cdot y = k \qquad y = \frac{k}{x} \]
        \textbf{Ejemplo:} 6 obreros terminan en 20 días. ¿8 obreros?
        \[ 6(20) = 8(x) \Rightarrow x = \destacar{15 \text{ días}} \]
      \end{alertblock}
    \end{column}
  \end{columns}

  \vspace{0.1cm}
  \begin{exampleblock}{Cómo decidir cuál es, sin equivocarse}
    \scriptsize
    Pregúntese: \destacar{si la primera magnitud aumenta, ¿la segunda aumenta o
    disminuye?} Más servidores $\Rightarrow$ más peticiones: \azul{directa}.
    Más obreros $\Rightarrow$ menos días: \azul{inversa}. \\
    {\tiny\color{uniGris} En la directa se multiplica en cruz; en la inversa,
    en línea recta.}
  \end{exampleblock}
\end{frame}

\begin{frame}{Regla de tres compuesta}
  \scriptsize
  \begin{exampleblock}{El problema}
    \tiny
    \destacar{4 programadores} desarrollan un módulo en \destacar{12 días},
    trabajando \destacar{6 horas diarias}. ¿En cuántos días lo harían
    \destacar{6 programadores} trabajando \destacar{8 horas diarias}?
  \end{exampleblock}

  \vspace{0.08cm}
  \begin{columns}[T,onlytextwidth]
    \begin{column}{0.48\textwidth}
      \begin{block}{Método del trabajo total}
        \tiny
        El trabajo no cambia; se mide en \azul{horas-persona}:
        \[ 4 \times 12 \times 6 = \destacar{288} \text{ h-persona} \]
        Con la nueva configuración:
        \[ 6 \times d \times 8 = 288 \;\Rightarrow\; 48d = 288 \]
        \[ d = \destacar{6 \text{ días}} \]
      \end{block}
    \end{column}
    \begin{column}{0.48\textwidth}
      \begin{alertblock}{Por qué este método y no las flechas}
        \tiny
        El método de las flechas exige decidir la dirección de cada magnitud y
        es donde todo el mundo se equivoca.
        \vspace{0.08cm}
        Calcular el \destacar{trabajo total} y volverlo a repartir es
        \azul{una sola ecuación}, sin flechas y sin dudas.
      \end{alertblock}
    \end{column}
  \end{columns}

  \vspace{0.08cm}
  {\tiny\color{uniGris} \destacar{Verificación de sentido:} más gente y más
  horas al día $\Rightarrow$ debe tardar \azul{menos} de 12 días. Salió 6:
  coherente. Si le hubiera dado 24, el resultado estaría mal aunque las cuentas
  parecieran bien.}
\end{frame}

\begin{frame}{Porcentajes: aumentos, descuentos y variación}
  \scriptsize
  \begin{columns}[T,onlytextwidth]
    \begin{column}{0.48\textwidth}
      \begin{block}{Aumentar y descontar en un paso}
        \tiny
        Descontar $15\%$ es multiplicar por $0{,}85$; \\
        aumentar $19\%$ es multiplicar por $1{,}19$.
        \vspace{0.06cm}
        Un equipo de \$850\,000 con 15\% de descuento:
        \[ 850\,000 (0{,}85) = \destacar{\$722\,500} \]
        Y si a eso se le suma el IVA del 19\%:
        \[ 722\,500(1{,}19) = \destacar{\$859\,775} \]
      \end{block}
    \end{column}
    \begin{column}{0.48\textwidth}
      \begin{alertblock}{Variación porcentual}
        \tiny
        \[ \frac{\text{final} - \text{inicial}}{\text{inicial}}\times 100 \]
        De 120 a 150: \ $\tfrac{30}{120}(100) = \destacar{+25\%}$

        \vspace{0.06cm}
        De 150 a 120: \ $\tfrac{-30}{150}(100) = \destacar{-20\%}$

        \vspace{0.06cm}
        \destacar{No son el mismo porcentaje}, porque la base cambió.
      \end{alertblock}
    \end{column}
  \end{columns}

  \vspace{0.08cm}
  \begin{exampleblock}{Descuentos sucesivos: la trampa comercial}
    \tiny
    ``20\% y además 10\%'' \destacar{no es} 30\%. Sobre \$100:
    $100(0{,}80)(0{,}90) = 72$, es decir un descuento real del
    \destacar{28\%}. \azul{Los porcentajes sucesivos se multiplican, no se
    suman.}
  \end{exampleblock}
\end{frame}

\begin{frame}{Repartos proporcionales y escalas}
  \scriptsize
  \begin{columns}[T,onlytextwidth]
    \begin{column}{0.50\textwidth}
      \begin{block}{Reparto proporcional}
        \tiny
        Repartir \$4\,500\,000 entre tres socios en razón $2:3:4$.
        \vspace{0.06cm}
        Partes totales: $2+3+4 = 9$
        \[ \text{una parte} = \frac{4\,500\,000}{9} = 500\,000 \]
        \[ \destacar{\$1\,000\,000},\quad \destacar{\$1\,500\,000},\quad
           \destacar{\$2\,000\,000} \]
        Verificación: suman \$4\,500\,000 \checkmark
      \end{block}
    \end{column}
    \begin{column}{0.46\textwidth}
      \begin{block}{Escalas}
        \tiny
        Una escala $1:200$ significa que 1 unidad en el plano equivale a 200 en
        la realidad.
        \vspace{0.06cm}
        Un muro mide 3,5 cm en el plano:
        \[ 3{,}5 \times 200 = 700 \text{ cm} = \destacar{7 \text{ m}} \]
      \end{block}
    \end{column}
  \end{columns}

  \vspace{0.08cm}
  \begin{alertblock}{Dónde aparece esto en su carrera}
    \tiny
    Los repartos proporcionales son la base del \destacar{reparto de recursos}
    entre procesos y del \destacar{balanceo de carga} entre servidores. Las
    escalas son las \destacar{transformaciones de coordenadas} de cualquier
    interfaz gráfica o motor de videojuego.
  \end{alertblock}
\end{frame}

\begin{frame}{Ejercicio 1 — en clase}
  \begin{alertblock}{En parejas (15 minutos)}
    \scriptsize
    \begin{enumerate}\setlength{\itemsep}{0.06cm}
      \item Halle $x$: \quad $\dfrac{x}{8} = \dfrac{21}{24}$
      \item Halle la media proporcional entre 4 y 9.
      \item Tres impresoras tardan 45 minutos en un lote. ¿Cuánto tardarían
            cinco impresoras iguales?
      \item Un carro recorre 45 km con 3 galones. ¿Cuántos km recorre con 8
            galones?
      \item 18 obreros terminan una obra en 24 días. Si solo llegan 12,
            ¿cuántos días tardarán?
      \item En una empresa la razón de desarrolladores a testers es $5:2$. Si
            hay 91 personas en total, ¿cuántos hay de cada uno?
    \end{enumerate}
  \end{alertblock}
\end{frame}

\begin{frame}{Ejercicio 1 — solución}
  \scriptsize
  \begin{columns}[T,onlytextwidth]
    \begin{column}{0.49\textwidth}
      \begin{block}{Puntos 1 a 3}
        \tiny
        \textbf{1.} $24x = 8(21) = 168 \Rightarrow x = \destacar{7}$

        \vspace{0.08cm}
        \textbf{2.} $x^2 = 4(9) = 36 \Rightarrow x = \destacar{6}$

        \vspace{0.08cm}
        \textbf{3.} \azul{Inversa}: más impresoras, menos tiempo.
        \[ 3(45) = 5x \Rightarrow x = \destacar{27 \text{ min}} \]
      \end{block}
    \end{column}
    \begin{column}{0.49\textwidth}
      \begin{block}{Puntos 4 a 6}
        \tiny
        \textbf{4.} \azul{Directa}: $\tfrac{45}{3} = \tfrac{x}{8}
        \Rightarrow x = \destacar{120 \text{ km}}$

        \vspace{0.08cm}
        \textbf{5.} \azul{Inversa}: $18(24) = 12x \Rightarrow
        x = \destacar{36 \text{ días}}$

        \vspace{0.08cm}
        \textbf{6.} $5+2 = 7$ partes; $91/7 = 13$.
        \[ \destacar{65} \text{ desarrolladores},\quad
           \destacar{26} \text{ testers} \]
      \end{block}
    \end{column}
  \end{columns}

  \vspace{0.08cm}
  \begin{alertblock}{Lo que decide cada punto}
    \tiny
    En el \destacar{3 y el 5} el error es plantear regla de tres directa. Antes
    de escribir la proporción, \azul{diga en voz alta si la segunda magnitud
    sube o baja}. Es un segundo de reflexión que salva el ejercicio.
  \end{alertblock}
\end{frame}

\begin{frame}{Ejercicio 2 — en clase}
  \begin{alertblock}{Individual (12 minutos)}
    \scriptsize
    \begin{enumerate}\setlength{\itemsep}{0.06cm}
      \item Un salario de \$2\,500\,000 sube el $8\%$. ¿Cuál es el nuevo
            salario?
      \item Un artículo de \$180\,000 tiene 25\% de descuento. ¿Cuánto se
            paga?
      \item Las ventas pasaron de 320 a 400 unidades. ¿Cuál fue la variación
            porcentual?
      \item Y si pasaran de 400 a 320, ¿cuál sería?
      \item Una tienda ofrece ``30\% y luego 20\% adicional''. ¿Cuál es el
            descuento real?
      \item Reparta \$7\,200\,000 entre tres áreas en razón $3:4:5$.
    \end{enumerate}
  \end{alertblock}
\end{frame}

\begin{frame}{Ejercicio 2 — solución}
  \scriptsize
  \begin{columns}[T,onlytextwidth]
    \begin{column}{0.49\textwidth}
      \begin{block}{Puntos 1 a 3}
        \tiny
        \textbf{1.} $2\,500\,000(1{,}08) = \destacar{\$2\,700\,000}$

        \vspace{0.08cm}
        \textbf{2.} $180\,000(0{,}75) = \destacar{\$135\,000}$

        \vspace{0.08cm}
        \textbf{3.} $\tfrac{400-320}{320}(100) = \destacar{+25\%}$
      \end{block}
    \end{column}
    \begin{column}{0.49\textwidth}
      \begin{block}{Puntos 4 a 6}
        \tiny
        \textbf{4.} $\tfrac{320-400}{400}(100) = \destacar{-20\%}$

        \vspace{0.08cm}
        \textbf{5.} $0{,}70 \times 0{,}80 = 0{,}56 \Rightarrow$ descuento del
        \destacar{44\%}, no del 50\%.

        \vspace{0.08cm}
        \textbf{6.} $3+4+5 = 12$; una parte $= 600\,000$:
        \[ \destacar{\$1{,}8\text{M}},\ \destacar{\$2{,}4\text{M}},\
           \destacar{\$3{,}0\text{M}} \]
      \end{block}
    \end{column}
  \end{columns}

  \vspace{0.08cm}
  \begin{alertblock}{Compare el 3 con el 4}
    \tiny
    Subir de 320 a 400 es $+25\%$, pero bajar de 400 a 320 es $-20\%$:
    \destacar{la misma diferencia de 80 unidades da porcentajes distintos},
    porque la base cambia. Por eso ``subió 25\% y luego bajó 25\%'' nunca
    devuelve al valor inicial.
  \end{alertblock}
\end{frame}

\seccion{Repaso integrador de la Unidad 1}

\begin{frame}{Mapa de la Unidad 1}
  \centering
  \begin{tikzpicture}[scale=0.92,every node/.style={font=\tiny,align=center},
    nodo/.style={draw=uniAzul,fill=uniAzul!10,rounded corners,inner sep=3pt,
                 minimum width=2.4cm,minimum height=0.62cm},
    flecha/.style={->,>=latex,uniGris}]
    \node[nodo] (a) at (0,0.9) {1.1 Conjuntos\\ numéricos};
    \node[nodo] (b) at (3.5,0.9) {1.2 Enteros};
    \node[nodo] (c) at (7.0,0.9) {1.3 Racionales};
    \node[nodo,fill=uniCremaClaro,draw=uniNaranja] (d) at (3.5,-0.9)
      {1.4 Conjuntos};
    \node[nodo,fill=uniCrema,draw=uniTitulo] (e) at (7.0,-0.9)
      {1.5 Razones y\\ proporciones};
    \node[nodo] (f) at (10.4,0) {Unidad 2:\\ potencias y\\ logaritmos};
    \draw[flecha] (a)--(b); \draw[flecha] (b)--(c);
    \draw[flecha] (a)--(d); \draw[flecha] (c)--(e);
    \draw[flecha] (c)--(f); \draw[flecha] (e)--(f);
  \end{tikzpicture}

  \vspace{0.25cm}
  \begin{alertblock}{La idea que atraviesa la unidad}
    \scriptsize
    Cada tema amplió lo que se podía hacer: de contar a restar
    ($\mathbb{Z}$), de restar a dividir ($\mathbb{Q}$), de operar números a
    operar colecciones (conjuntos) y de operar a \destacar{comparar}
    (proporciones). \azul{La unidad 2 amplía otra vez: multiplicar muchas
    veces.}
  \end{alertblock}
\end{frame}

\begin{frame}{Formulario de la Unidad 1}
  \tiny
  \begin{columns}[T,onlytextwidth]
    \begin{column}{0.48\textwidth}
      \begin{block}{Números y operaciones}
        \tiny
        \[ \mathbb{N}\subset\mathbb{Z}\subset\mathbb{Q}\subset\mathbb{R},
           \qquad \mathbb{R}=\mathbb{Q}\cup\mathbb{I} \]
        \[ a-b = a+(-b) \qquad -a^{n} \neq (-a)^{n} \ (n \text{ par}) \]
        \[ \frac{a}{b}\pm\frac{c}{d} = \frac{ad\pm bc}{bd} \qquad
           \frac{a}{b}\div\frac{c}{d} = \frac{a}{b}\cdot\frac{d}{c} \]
        Jerarquía: $(\ )$, potencias, $\times\div$, $+-$, \azul{izquierda a
        derecha}.
      \end{block}
    \end{column}
    \begin{column}{0.48\textwidth}
      \begin{block}{Conjuntos y proporciones}
        \tiny
        \[ |\mathcal{P}(A)| = 2^{|A|} \qquad
           (A\cup B)^{c} = A^{c}\cap B^{c} \]
        \[ |A\cup B| = |A|+|B|-|A\cap B| \]
        \[ \frac{a}{b}=\frac{c}{d} \Leftrightarrow ad = bc \]
        Directa: $\tfrac{y}{x}=k$. \quad Inversa: $xy=k$.
        \[ \text{Var.\%} = \frac{\text{final}-\text{inicial}}{\text{inicial}}
           \times 100 \]
      \end{block}
    \end{column}
  \end{columns}

  \vspace{0.05cm}
  {\tiny\color{uniGris} MCD: factores comunes con \azul{menor} exponente. \quad
  mcm: todos los factores con \azul{mayor} exponente.}
\end{frame}

\begin{frame}{Simulacro del Examen 1 — parte A}
  \begin{alertblock}{Individual (20 minutos, sin apuntes)}
    \scriptsize
    \begin{enumerate}\setlength{\itemsep}{0.02cm}
      \item Clasifique en $\mathbb{N}$, $\mathbb{Z}$, $\mathbb{Q}$,
            $\mathbb{I}$: \ $-\sqrt{49}$, \ $\sqrt{5}$, \ $\tfrac{27}{9}$,
            \ $0{,}\overline{81}$.
      \item Convierta $0{,}\overline{81}$ a fracción.
      \item Calcule: \ $-3^{2} + (-2)^{3} - 12\div(-4)\cdot 2$
      \item Calcule: \ $\dfrac{5}{6} - \dfrac{3}{4} \cdot \dfrac{2}{9}$
      \item Calcule: \ $\dfrac{\ \frac{2}{3}+\frac{1}{6}\ }
            {\ \frac{5}{9}\ }$
      \item Halle $\mathrm{MCD}(42,70)$ y $\mathrm{mcm}(42,70)$.
    \end{enumerate}
  \end{alertblock}
\end{frame}

\begin{frame}{Simulacro — solución de la parte A}
  \scriptsize
  \begin{columns}[T,onlytextwidth]
    \begin{column}{0.49\textwidth}
      \begin{block}{Puntos 1 a 3}
        \tiny
        \textbf{1.} $-\sqrt{49}=-7$: $\mathbb{Z},\mathbb{Q}$. \
        $\sqrt{5}$: $\mathbb{I}$. \ $\tfrac{27}{9}=3$:
        $\mathbb{N},\mathbb{Z},\mathbb{Q}$. \ $0{,}\overline{81}$:
        $\mathbb{Q}$.

        \vspace{0.08cm}
        \textbf{2.} $100x - x = 81 \Rightarrow x = \tfrac{81}{99}
        = \destacar{\tfrac{9}{11}}$

        \vspace{0.08cm}
        \textbf{3.} $-9 + (-8) - (-3)(2) = -9-8+6 = \destacar{-11}$
      \end{block}
    \end{column}
    \begin{column}{0.49\textwidth}
      \begin{block}{Puntos 4 a 6}
        \tiny
        \textbf{4.} $\tfrac{5}{6} - \tfrac{6}{36}
        = \tfrac{5}{6}-\tfrac{1}{6} = \destacar{\tfrac{2}{3}}$

        \vspace{0.08cm}
        \textbf{5.} $\tfrac{5}{6}\div\tfrac{5}{9}
        = \tfrac{5}{6}\cdot\tfrac{9}{5} = \destacar{\tfrac{3}{2}}$

        \vspace{0.08cm}
        \textbf{6.} $42 = 2\cdot 3\cdot 7$, \ $70 = 2\cdot 5\cdot 7$
        \[ \mathrm{MCD} = 14 \qquad \mathrm{mcm} = 210 \]
      \end{block}
    \end{column}
  \end{columns}

  \vspace{0.05cm}
  {\tiny\color{uniGris} En el \destacar{3}, $12\div(-4)\cdot 2 = -3\cdot 2
  = -6$, y como va restando queda $+6$. Verificación del 6:
  $14\cdot 210 = 2\,940 = 42\cdot 70$ \checkmark}
\end{frame}

\begin{frame}{Simulacro del Examen 1 — parte B}
  \begin{alertblock}{Individual (20 minutos)}
    \scriptsize
    \begin{enumerate}\setlength{\itemsep}{0.02cm}
      \item Con $U=\{1,\dots,10\}$, $A=\{2,3,5,7\}$ y $B=\{1,2,3,4\}$,
            determine $A\cup B$, $A\cap B$, $A-B$ y $(A\cup B)^{c}$.
      \item Verifique en ese caso una de las leyes de De Morgan.
      \item De 90 estudiantes, 55 usan Git, 40 usan Docker y 20 usan ambos.
            ¿Cuántos no usan ninguno? ¿Cuántos usan solo Git?
      \item Cinco máquinas producen 800 piezas en un turno. ¿Cuántas producen
            ocho máquinas iguales?
      \item Nueve operarios terminan un pedido en 12 días. ¿Cuántos días
            tardarían seis operarios?
      \item Un producto de \$240\,000 sube 15\% y luego baja 15\%.
            ¿Queda en \$240\,000? Justifique con el cálculo.
    \end{enumerate}
  \end{alertblock}
\end{frame}

\begin{frame}{Simulacro — solución de la parte B}
  \scriptsize
  \begin{columns}[T,onlytextwidth]
    \begin{column}{0.49\textwidth}
      \begin{block}{Puntos 1 a 3}
        \tiny
        \textbf{1.} $A\cup B = \{1,2,3,4,5,7\}$; \ $A\cap B = \{2,3\}$; \
        $A-B = \{5,7\}$; \ $(A\cup B)^c = \{6,8,9,10\}$

        \vspace{0.08cm}
        \textbf{2.} $A^c\cap B^c = \{1,4,6,8,9,10\}\cap\{5,6,7,8,9,10\}
        = \{6,8,9,10\}$ \checkmark

        \vspace{0.08cm}
        \textbf{3.} $|G\cup D| = 55+40-20 = 75$; ninguno:
        $90-75 = \destacar{15}$; solo Git: $55-20 = \destacar{35}$
      \end{block}
    \end{column}
    \begin{column}{0.49\textwidth}
      \begin{block}{Puntos 4 a 6}
        \tiny
        \textbf{4.} Directa: $\tfrac{800}{5} = \tfrac{x}{8}
        \Rightarrow \destacar{1\,280}$ piezas

        \vspace{0.08cm}
        \textbf{5.} Inversa: $9(12) = 6x \Rightarrow \destacar{18}$ días

        \vspace{0.08cm}
        \textbf{6.} $240\,000(1{,}15)(0{,}85) = 234\,600$:
        \destacar{no vuelve al inicial}, queda \$5\,400 por debajo, porque el
        15\% de bajada se calcula sobre una base mayor.
      \end{block}
    \end{column}
  \end{columns}

  \vspace{0.05cm}
  {\tiny\color{uniGris} El punto 6 es el que más se falla del examen, y no por
  cuentas: por \destacar{creer que los porcentajes se cancelan}.}
\end{frame}

\begin{frame}{Recomendaciones para el Examen 1}
  \begin{block}{Qué se evalúa}
    \tiny
    \begin{itemize}\setlength{\itemsep}{-0.02cm}
      \item Clasificar números y convertir decimales periódicos a fracción.
      \item Operar enteros y racionales respetando signos y jerarquía.
      \item Calcular MCD y mcm y usarlos en problemas.
      \item Operar conjuntos, aplicar De Morgan y contar con Venn.
      \item Plantear proporciones, regla de tres directa e inversa y
            porcentajes.
    \end{itemize}
  \end{block}

  \vspace{0.08cm}
  \begin{alertblock}{Seis consejos concretos}
    \tiny
    (1) Escriba \destacar{un paso por línea}: los errores de signo aparecen al
    saltarse pasos. \quad
    (2) En jerarquía, respete \azul{izquierda a derecha}. \quad
    (3) Simplifique siempre el resultado final. \quad
    (4) En proporciones, pregúntese primero si es directa o inversa. \quad
    (5) \destacar{Verifique}: sustituya la respuesta en el enunciado. \quad
    (6) Si un resultado no tiene sentido (más obreros y más días), revise antes
    de entregar.
  \end{alertblock}

  \vspace{0.05cm}
  {\tiny\color{uniGris} Examen 1: semana 6 (14 al 19 de septiembre), 35\% de
  la nota, en Aulas Virtuales (Moodle), presencial en el aula.}
\end{frame}

\begin{frame}{Ejercicios propuestos}
  \vspace{-0.1cm}
  \begin{beamercolorbox}[wd=\textwidth,sep=0.10cm]{block body}
    \centering\scriptsize
    \textbf{1--7:} proporciones. \quad \textbf{8--14:} regla de tres. \quad
    \textbf{15--20:} porcentajes y repaso de la unidad.
  \end{beamercolorbox}
  \listatiny
  \vspace{0.15cm}
  \begin{columns}[T,onlytextwidth]
    \begin{column}{0.47\textwidth}
      \begin{enumerate}\setlength{\itemsep}{0.16cm}
        \item Simplifique la razón $15:25$
        \item Halle $x$: $\dfrac{3}{7}=\dfrac{x}{28}$
        \item Halle $x$: $\dfrac{x}{5}=\dfrac{12}{20}$
        \item Media proporcional entre 9 y 25
        \item Reparta 320 en razón $3:5$
        \item Reparta \$6\,000\,000 en razón $1:2:3$
        \item Si $a:b = 4:7$ y $b = 35$, halle $a$
        \item 4 kg cuestan \$18\,000. ¿Y 7 kg?
        \item 6 llaves llenan un tanque en 8 h. ¿Y 4 llaves?
        \item 3 buses llevan 120 pasajeros. ¿Y 7 buses?
      \end{enumerate}
    \end{column}
    \begin{column}{0.50\textwidth}
      \begin{enumerate}\setlength{\itemsep}{0.16cm}
        \setcounter{enumi}{10}
        \item 15 obreros en 8 días. ¿Cuántos días 10 obreros?
        \item 5 programadores, 20 días, 6 h/día. ¿Cuántos días con 8
              programadores a 5 h/día?
        \item Escala $1:50$; en el plano 4,2 cm. ¿Cuánto en la realidad?
        \item Un mapa $1:25\,000$: 6 cm equivalen a\dots
        \item El 12\% de 4\,500
        \item ¿Qué porcentaje es 45 de 300?
        \item De 250 a 300: variación porcentual
        \item De 300 a 250: variación porcentual
        \item Descuentos sucesivos de 10\% y 25\%: ¿descuento real?
        \item Un precio con IVA (19\%) es \$595\,000. ¿Cuál era sin IVA?
      \end{enumerate}
    \end{column}
  \end{columns}
\end{frame}

\begin{frame}{Respuestas}{Autoevaluación de los ejercicios propuestos}
  \listatiny
  \vspace{0.15cm}
  \begin{columns}[T,onlytextwidth]
    \begin{column}{0.47\textwidth}
      \begin{enumerate}\setlength{\itemsep}{0.16cm}
        \item $3:5$
        \item $x = 12$
        \item $x = 3$
        \item $\sqrt{225} = 15$
        \item $120$ y $200$
        \item \$1M, \$2M y \$3M
        \item $a = 20$
        \item $\$31\,500$
        \item $12$ h \ (inversa)
        \item $280$ pasajeros
      \end{enumerate}
    \end{column}
    \begin{column}{0.50\textwidth}
      \begin{enumerate}\setlength{\itemsep}{0.16cm}
        \setcounter{enumi}{10}
        \item $12$ días
        \item $5(20)(6)=600$; \ $600/(8\cdot 5) = 15$ días
        \item $4{,}2\times 50 = 210$ cm $= 2{,}1$ m
        \item $6\times 25\,000 = 150\,000$ cm $= 1{,}5$ km
        \item $540$
        \item $15\%$
        \item $+20\%$
        \item $-16{,}\overline{6}\%$
        \item $0{,}90\times 0{,}75 = 0{,}675 \Rightarrow 32{,}5\%$
        \item $595\,000 / 1{,}19 = \$500\,000$
      \end{enumerate}
    \end{column}
  \end{columns}
\end{frame}

\begin{frame}{Referencias}
  \footnotesize
  \begin{itemize}
    \item[{[1]}] Ramos del Olmo, Á. M. (2016). \textit{Matemáticas básicas
      para el acceso a la universidad}. Ediciones Pirámide.
      \texttt{https://elibro.net/es/ereader/uniminuto/49134}
    \item[{[2]}] González Pulido, J. W., Chacón Benavides, J. A. \&
      Fonseca Correa, L. Á. (2023). \textit{Matemática básica} (1.\textsuperscript{a} ed.).
      Universidad Mariana.
      \texttt{https://elibro.net/es/ereader/uniminuto/280647}
    \item[{[3]}] Placencia Valero, J. (2008). \textit{Compendio de matemática
      básica elemental} (1.\textsuperscript{a} ed.). Editorial Tébar Flores.
      \texttt{https://elibro.net/es/ereader/uniminuto/269758}
  \end{itemize}

  \vspace{0.35cm}
  \begin{center}
    {\Large\color{uniAzul}\textbf{Éxitos en el Examen 1}}
  \end{center}
\end{frame}

\end{document}
